\documentclass[11pt]{article}
% The preamble every project's paper shares: % The preamble every project's paper shares: % The preamble every project's paper shares: \input{../../preamble} from
% projects/<project>/paper/. Boards come from chessboard and xskak, plots from
% pgfplots; every number a paper prints comes from its generated.tex.
\usepackage[LSBC4,T1]{fontenc}
\usepackage[utf8]{inputenc}
\usepackage{lmodern}
\usepackage[margin=1in]{geometry}
\usepackage{microtype}
\usepackage{amsmath}
\usepackage{booktabs}
\usepackage{tabularx}
\usepackage{array}
\usepackage{xcolor}
\usepackage[obeyspaces,spaces]{url}
\usepackage{xskak}
\usepackage{pgfplots}
\pgfplotsset{compat=1.18}
\usepackage{caption}
\usepackage{subcaption}
\usepackage{enumitem}
\usepackage[section]{placeins}
\usepackage[hidelinks]{hyperref}

\setchessboard{boardfontencoding=LSBC4, boardfontfamily=skaknew, blackfieldcolor=black!22,
  whitefieldcolor=white, setfontcolors, showmover=false, label=true, labelfontsize=6pt,
  boardfontsize=14pt}

% A move or line in UCI, and a position in FEN, both breakable at a slash.
\newcommand{\uci}[1]{\texttt{#1}}
\newcommand{\fen}[1]{{\footnotesize\path{#1}}}
\newcolumntype{L}{>{\raggedright\arraybackslash}X}

\setlength{\parskip}{0.4em}
\setlength{\parindent}{0pt}
 from
% projects/<project>/paper/. Boards come from chessboard and xskak, plots from
% pgfplots; every number a paper prints comes from its generated.tex.
\usepackage[LSBC4,T1]{fontenc}
\usepackage[utf8]{inputenc}
\usepackage{lmodern}
\usepackage[margin=1in]{geometry}
\usepackage{microtype}
\usepackage{amsmath}
\usepackage{booktabs}
\usepackage{tabularx}
\usepackage{array}
\usepackage{xcolor}
\usepackage[obeyspaces,spaces]{url}
\usepackage{xskak}
\usepackage{pgfplots}
\pgfplotsset{compat=1.18}
\usepackage{caption}
\usepackage{subcaption}
\usepackage{enumitem}
\usepackage[section]{placeins}
\usepackage[hidelinks]{hyperref}

\setchessboard{boardfontencoding=LSBC4, boardfontfamily=skaknew, blackfieldcolor=black!22,
  whitefieldcolor=white, setfontcolors, showmover=false, label=true, labelfontsize=6pt,
  boardfontsize=14pt}

% A move or line in UCI, and a position in FEN, both breakable at a slash.
\newcommand{\uci}[1]{\texttt{#1}}
\newcommand{\fen}[1]{{\footnotesize\path{#1}}}
\newcolumntype{L}{>{\raggedright\arraybackslash}X}

\setlength{\parskip}{0.4em}
\setlength{\parindent}{0pt}
 from
% projects/<project>/paper/. Boards come from chessboard and xskak, plots from
% pgfplots; every number a paper prints comes from its generated.tex.
\usepackage[LSBC4,T1]{fontenc}
\usepackage[utf8]{inputenc}
\usepackage{lmodern}
\usepackage[margin=1in]{geometry}
\usepackage{microtype}
\usepackage{amsmath}
\usepackage{booktabs}
\usepackage{tabularx}
\usepackage{array}
\usepackage{xcolor}
\usepackage[obeyspaces,spaces]{url}
\usepackage{xskak}
\usepackage{pgfplots}
\pgfplotsset{compat=1.18}
\usepackage{caption}
\usepackage{subcaption}
\usepackage{enumitem}
\usepackage[section]{placeins}
\usepackage[hidelinks]{hyperref}

\setchessboard{boardfontencoding=LSBC4, boardfontfamily=skaknew, blackfieldcolor=black!22,
  whitefieldcolor=white, setfontcolors, showmover=false, label=true, labelfontsize=6pt,
  boardfontsize=14pt}

% A move or line in UCI, and a position in FEN, both breakable at a slash.
\newcommand{\uci}[1]{\texttt{#1}}
\newcommand{\fen}[1]{{\footnotesize\path{#1}}}
\newcolumntype{L}{>{\raggedright\arraybackslash}X}

\setlength{\parskip}{0.4em}
\setlength{\parindent}{0pt}

% Written by projects/think-time-rapid/paper/generate.mjs; do not edit by hand.

\newcommand{\AppCommit}{d8b771238c}
\newcommand{\Month}{2026-09}
\newcommand{\WindowLo}{1250}
\newcommand{\WindowHi}{1449}
\newcommand{\MinClock}{60}
\newcommand{\Strata}{five}
\newcommand{\Draws}{2,000}
\newcommand{\Seed}{20261014}
\newcommand{\Plateau}{60}
\newcommand{\BarShare}{80}
\newcommand{\DefendingPawns}{3}
\newcommand{\EarlyOwn}{four}
\newcommand{\WindowCentre}{1350}
\newcommand{\RatingLo}{900}
\newcommand{\RatingHi}{2099}
\newcommand{\Rushed}{10}
\newcommand{\Deep}{30}
\newcommand{\BookEnd}{45}
\newcommand{\Decided}{5}
\newcommand{\GamesRead}{6,185,560}
\newcommand{\GamesRapid}{1,019,414}
\newcommand{\GamesWithEvals}{195,409}
\newcommand{\GamesKept}{74,000}
\newcommand{\GamesBuilt}{74,000}
\newcommand{\BlackChecked}{800}
\newcommand{\MovesTimed}{4,721,351}
\newcommand{\MovesAnalysed}{3,738,138}
\newcommand{\MovesPrimary}{980,946}
\newcommand{\PlayerGames}{39,902}
\newcommand{\PrimaryQuick}{668,544}
\newcommand{\PrimaryLong}{54,331}
\newcommand{\LongShare}{6}
\newcommand{\QuickShare}{68}
\newcommand{\KeptTable}{%
\begin{tabular}{lrrrrrrr}
\toprule
mean rating & 800--999 & 1000--1199 & 1200--1399 & 1400--1599 & 1600--1799 & 1800--1999 & 2000--2199 \\
\midrule
games & 6,000 & 8,000 & 20,000 & 20,000 & 8,000 & 6,000 & 6,000 \\
\bottomrule
\end{tabular}
}
\newcommand{\AucQuick}{0.799}
\newcommand{\AucOther}{0.762}
\newcommand{\Trees}{400}
\newcommand{\SeriousQuick}{10.9}
\newcommand{\TriggerTable}{%
\begin{tabular}{lrrrrrllr}
\toprule
 & per & \multicolumn{2}{c}{moves} & \multicolumn{2}{c}{serious errors} & & & Holm \\
\cmidrule(lr){3-4}\cmidrule(lr){5-6}
moment & game & quick & long & quick & long & $\Delta$ & $\Delta - \Delta_{\text{none}}$ & $p$ \\
\midrule
Book end & 1.0 & 33,996 & 437 & 5.0 & 12.1 & $+$1.6 {\footnotesize ($-$0.1, $+$3.4)} & $-$3.1 {\footnotesize ($-$4.9, $-$1.2)} & 0.008 \\
Conversion & 5.3 & 134,722 & 13,192 & 11.9 & 23.4 & $+$8.8 {\footnotesize ($+$8.2, $+$9.5)} & $+$4.1 {\footnotesize ($+$3.3, $+$5.0)} & 0.004 \\
Forcing & 11.2 & 306,030 & 24,636 & 13.0 & 26.9 & $+$7.6 {\footnotesize ($+$7.1, $+$8.0)} & $+$2.9 {\footnotesize ($+$2.1, $+$3.6)} & 0.004 \\
Structure & 0.2 & 4,871 & 307 & 9.7 & 13.7 & $+$0.6 {\footnotesize ($-$2.3, $+$3.7)} & $-$4.1 {\footnotesize ($-$6.9, $-$1.0)} & 0.010 \\
\midrule
none & 7.0 & 188,925 & 15,759 & 11.0 & 23.1 & $+$4.7 {\footnotesize ($+$4.1, $+$5.3)} &  &  \\
\bottomrule
\end{tabular}
}
\newcommand{\BookEndDelta}{$+$1.6}
\newcommand{\BookEndLo}{$-$0.1}
\newcommand{\BookEndHi}{$+$3.4}
\newcommand{\BookEndVsNone}{$-$3.1}
\newcommand{\BookEndVsNoneLo}{$-$4.9}
\newcommand{\BookEndVsNoneHi}{$-$1.2}
\newcommand{\BookEndPHolm}{0.008}
\newcommand{\BookEndPerGame}{1.0}
\newcommand{\BookEndVerdict}{no sign it pays}
\newcommand{\BookEndThinkGap}{44}
\newcommand{\BookEndRateQuick}{5.0}
\newcommand{\BookEndRateLong}{12.1}
\newcommand{\BookEndDropped}{0.0}
\newcommand{\BookEndSaved}{$-$0.021}
\newcommand{\BookEndSavedLo}{$-$0.047}
\newcommand{\BookEndSavedHi}{0.002}
\newcommand{\ConversionDelta}{$+$8.8}
\newcommand{\ConversionLo}{$+$8.2}
\newcommand{\ConversionHi}{$+$9.5}
\newcommand{\ConversionVsNone}{$+$4.1}
\newcommand{\ConversionVsNoneLo}{$+$3.3}
\newcommand{\ConversionVsNoneHi}{$+$5.0}
\newcommand{\ConversionPHolm}{0.004}
\newcommand{\ConversionPerGame}{5.3}
\newcommand{\ConversionVerdict}{no sign it pays}
\newcommand{\ConversionThinkGap}{41}
\newcommand{\ConversionRateQuick}{11.9}
\newcommand{\ConversionRateLong}{23.4}
\newcommand{\ConversionDropped}{0.0}
\newcommand{\ConversionSaved}{$-$0.128}
\newcommand{\ConversionSavedLo}{$-$0.137}
\newcommand{\ConversionSavedHi}{$-$0.118}
\newcommand{\ForcingDelta}{$+$7.6}
\newcommand{\ForcingLo}{$+$7.1}
\newcommand{\ForcingHi}{$+$8.0}
\newcommand{\ForcingVsNone}{$+$2.9}
\newcommand{\ForcingVsNoneLo}{$+$2.1}
\newcommand{\ForcingVsNoneHi}{$+$3.6}
\newcommand{\ForcingPHolm}{0.004}
\newcommand{\ForcingPerGame}{11.2}
\newcommand{\ForcingVerdict}{no sign it pays}
\newcommand{\ForcingThinkGap}{43}
\newcommand{\ForcingRateQuick}{13.0}
\newcommand{\ForcingRateLong}{26.9}
\newcommand{\ForcingDropped}{0.0}
\newcommand{\ForcingSaved}{$-$0.106}
\newcommand{\ForcingSavedLo}{$-$0.113}
\newcommand{\ForcingSavedHi}{$-$0.100}
\newcommand{\StructureDelta}{$+$0.6}
\newcommand{\StructureLo}{$-$2.3}
\newcommand{\StructureHi}{$+$3.7}
\newcommand{\StructureVsNone}{$-$4.1}
\newcommand{\StructureVsNoneLo}{$-$6.9}
\newcommand{\StructureVsNoneHi}{$-$1.0}
\newcommand{\StructurePHolm}{0.010}
\newcommand{\StructurePerGame}{0.2}
\newcommand{\StructureVerdict}{no sign it pays}
\newcommand{\StructureThinkGap}{44}
\newcommand{\StructureRateQuick}{9.7}
\newcommand{\StructureRateLong}{13.7}
\newcommand{\StructureDropped}{0.0}
\newcommand{\StructureSaved}{$-$0.008}
\newcommand{\StructureSavedLo}{$-$0.051}
\newcommand{\StructureSavedHi}{0.031}
\newcommand{\NoneDelta}{$+$4.7}
\newcommand{\NoneLo}{$+$4.1}
\newcommand{\NoneHi}{$+$5.3}
\newcommand{\NonePerGame}{7.0}
\newcommand{\NoneRateQuick}{11.0}
\newcommand{\NoneRateLong}{23.1}
\newcommand{\NoneThinkGap}{41}
\newcommand{\Kept}{zero}
\newcommand{\NoSign}{four}
\newcommand{\CandidateTable}{%
\begin{tabular}{lrrrllrc}
\toprule
 & per & \multicolumn{2}{c}{moves} & & & Holm & \\
\cmidrule(lr){3-4}
before the move & game & quick & long & $\Delta$ & $\Delta - \Delta_{\text{none}}$ & $p$ & joins \\
\midrule
in check & 1.5 & 50,679 & 959 & $+$3.7 {\footnotesize ($+$1.7, $+$5.7)} & $-$1.0 {\footnotesize ($-$3.1, $+$1.0)} & 0.648 & no \\
after a capture & 6.3 & 194,911 & 9,034 & $+$8.2 {\footnotesize ($+$7.5, $+$8.9)} & $+$3.5 {\footnotesize ($+$2.6, $+$4.4)} & 0.008 & no \\
his piece hanging & 9.4 & 241,606 & 25,008 & $+$8.7 {\footnotesize ($+$8.2, $+$9.2)} & $+$4.0 {\footnotesize ($+$3.3, $+$4.7)} & 0.008 & no \\
opponent's piece hanging & 6.3 & 186,669 & 11,791 & $+$10.6 {\footnotesize ($+$9.9, $+$11.3)} & $+$5.9 {\footnotesize ($+$5.0, $+$6.8)} & 0.008 & no \\
after the opponent's error & 3.4 & 85,675 & 8,489 & $+$16.3 {\footnotesize ($+$15.1, $+$17.5)} & $+$11.6 {\footnotesize ($+$10.4, $+$12.8)} & 0.008 & no \\
defending, $-$3 or worse & 4.5 & 116,842 & 11,270 & $+$2.3 {\footnotesize ($+$1.9, $+$2.7)} & $-$2.4 {\footnotesize ($-$3.0, $-$1.8)} & 0.008 & no \\
book end and his next four & 4.8 & 159,050 & 4,040 & $+$4.6 {\footnotesize ($+$3.7, $+$5.5)} & $-$0.1 {\footnotesize ($-$1.0, $+$0.8)} & 0.901 & no \\
high danger, top fifth & 5.7 & 133,709 & 18,266 & $+$12.9 {\footnotesize ($+$12.2, $+$13.7)} & $+$8.2 {\footnotesize ($+$7.4, $+$9.1)} & 0.008 & no \\
\bottomrule
\end{tabular}
}
\newcommand{\CandidatesJoin}{zero}
\newcommand{\CandidatesPay}{zero}
\newcommand{\CandInCheckDelta}{$+$3.7}
\newcommand{\CandInCheckLo}{$+$1.7}
\newcommand{\CandInCheckHi}{$+$5.7}
\newcommand{\CandInCheckVsNone}{$-$1.0}
\newcommand{\CandInCheckPHolm}{0.648}
\newcommand{\CandAfterCaptureDelta}{$+$8.2}
\newcommand{\CandAfterCaptureLo}{$+$7.5}
\newcommand{\CandAfterCaptureHi}{$+$8.9}
\newcommand{\CandAfterCaptureVsNone}{$+$3.5}
\newcommand{\CandAfterCapturePHolm}{0.008}
\newcommand{\CandOwnHangingDelta}{$+$8.7}
\newcommand{\CandOwnHangingLo}{$+$8.2}
\newcommand{\CandOwnHangingHi}{$+$9.2}
\newcommand{\CandOwnHangingVsNone}{$+$4.0}
\newcommand{\CandOwnHangingPHolm}{0.008}
\newcommand{\CandOppHangingDelta}{$+$10.6}
\newcommand{\CandOppHangingLo}{$+$9.9}
\newcommand{\CandOppHangingHi}{$+$11.3}
\newcommand{\CandOppHangingVsNone}{$+$5.9}
\newcommand{\CandOppHangingPHolm}{0.008}
\newcommand{\CandAfterOppErrorDelta}{$+$16.3}
\newcommand{\CandAfterOppErrorLo}{$+$15.1}
\newcommand{\CandAfterOppErrorHi}{$+$17.5}
\newcommand{\CandAfterOppErrorVsNone}{$+$11.6}
\newcommand{\CandAfterOppErrorPHolm}{0.008}
\newcommand{\CandDefendingDelta}{$+$2.3}
\newcommand{\CandDefendingLo}{$+$1.9}
\newcommand{\CandDefendingHi}{$+$2.7}
\newcommand{\CandDefendingVsNone}{$-$2.4}
\newcommand{\CandDefendingPHolm}{0.008}
\newcommand{\CandEarlyDelta}{$+$4.6}
\newcommand{\CandEarlyLo}{$+$3.7}
\newcommand{\CandEarlyHi}{$+$5.5}
\newcommand{\CandEarlyVsNone}{$-$0.1}
\newcommand{\CandEarlyPHolm}{0.901}
\newcommand{\CandHighDangerDelta}{$+$12.9}
\newcommand{\CandHighDangerLo}{$+$12.2}
\newcommand{\CandHighDangerHi}{$+$13.7}
\newcommand{\CandHighDangerVsNone}{$+$8.2}
\newcommand{\CandHighDangerPHolm}{0.008}
\newcommand{\HighDangerCut}{20.2}
\newcommand{\BookEndBarPlot}{%
\begin{tikzpicture}
\begin{axis}[scale only axis, width=0.8\linewidth, height=3.6cm, title={Book end}, xtick={0,...,6}, xticklabels={10--14,15--19,20--29,30--44,45--59,60--89,90+},
  xmin=-0.5, xmax=6.5, x tick label style={font=\scriptsize}, ylabel={$\Delta$, points}, xlabel={think, seconds}, grid=major]
\addplot[only marks, mark=*, mark size=1.5pt, black, error bars/.cd, y dir=both, y explicit] coordinates {
(0, 2.34) += (0, 0.87) -= (0, 0.89)
(1, 3.85) += (0, 1.55) -= (0, 1.55)
(2, 2.94) += (0, 1.49) -= (0, 1.44)
(3, 1.84) += (0, 2.36) -= (0, 2.17)
(4, 2.09) += (0, 5.72) -= (0, 4.43)
(5, 0.03) += (0, 4.78) -= (0, 4.03)
(6, 0.09) += (0, 9.52) -= (0, 6.44)
};
\draw[gray, densely dotted] (axis cs:-1,0) -- (axis cs:7,0);
\end{axis}
\end{tikzpicture}
}
\newcommand{\ConversionBarPlot}{%
\begin{tikzpicture}
\begin{axis}[scale only axis, width=0.8\linewidth, height=3.6cm, title={Conversion}, xtick={0,...,6}, xticklabels={10--14,15--19,20--29,30--44,45--59,60--89,90+},
  xmin=-0.5, xmax=6.5, x tick label style={font=\scriptsize}, ylabel={$\Delta$, points}, xlabel={think, seconds}, grid=major]
\addplot[only marks, mark=*, mark size=1.5pt, black, error bars/.cd, y dir=both, y explicit] coordinates {
(0, 3.67) += (0, 0.40) -= (0, 0.41)
(1, 5.63) += (0, 0.56) -= (0, 0.54)
(2, 6.37) += (0, 0.56) -= (0, 0.58)
(3, 7.85) += (0, 0.77) -= (0, 0.74)
(4, 8.56) += (0, 1.40) -= (0, 1.35)
(5, 13.60) += (0, 2.11) -= (0, 2.09)
(6, 11.91) += (0, 3.51) -= (0, 3.41)
};
\draw[gray, densely dotted] (axis cs:-1,0) -- (axis cs:7,0);
\end{axis}
\end{tikzpicture}
}
\newcommand{\ForcingBarPlot}{%
\begin{tikzpicture}
\begin{axis}[scale only axis, width=0.8\linewidth, height=3.6cm, title={Forcing}, xtick={0,...,6}, xticklabels={10--14,15--19,20--29,30--44,45--59,60--89,90+},
  xmin=-0.5, xmax=6.5, x tick label style={font=\scriptsize}, ylabel={$\Delta$, points}, xlabel={think, seconds}, grid=major]
\addplot[only marks, mark=*, mark size=1.5pt, black, error bars/.cd, y dir=both, y explicit] coordinates {
(0, 3.52) += (0, 0.29) -= (0, 0.30)
(1, 4.54) += (0, 0.40) -= (0, 0.37)
(2, 5.40) += (0, 0.42) -= (0, 0.43)
(3, 7.14) += (0, 0.60) -= (0, 0.60)
(4, 8.50) += (0, 0.98) -= (0, 1.04)
(5, 7.42) += (0, 1.28) -= (0, 1.22)
(6, 9.96) += (0, 2.33) -= (0, 2.20)
};
\draw[gray, densely dotted] (axis cs:-1,0) -- (axis cs:7,0);
\end{axis}
\end{tikzpicture}
}
\newcommand{\StructureBarPlot}{%
\begin{tikzpicture}
\begin{axis}[scale only axis, width=0.8\linewidth, height=3.6cm, title={Structure}, xtick={0,...,6}, xticklabels={10--14,15--19,20--29,30--44,45--59,60--89,90+},
  xmin=-0.5, xmax=6.5, x tick label style={font=\scriptsize}, ylabel={$\Delta$, points}, xlabel={think, seconds}, grid=major]
\addplot[only marks, mark=*, mark size=1.5pt, black, error bars/.cd, y dir=both, y explicit] coordinates {
(0, 2.41) += (0, 2.26) -= (0, 2.18)
(1, 3.51) += (0, 2.98) -= (0, 2.66)
(2, 0.81) += (0, 2.89) -= (0, 2.72)
(3, 0.83) += (0, 3.96) -= (0, 3.59)
(4, 1.06) += (0, 7.35) -= (0, 6.82)
(5, -1.12) += (0, 6.54) -= (0, 5.87)
(6, 0.32) += (0, 13.15) -= (0, 12.00)
};
\draw[gray, densely dotted] (axis cs:-1,0) -- (axis cs:7,0);
\end{axis}
\end{tikzpicture}
}
\newcommand{\BookEndPlateau}{$+$0.1}
\newcommand{\BookEndPlateauLo}{$-$3.4}
\newcommand{\BookEndPlateauHi}{$+$3.9}
\newcommand{\BookEndBar}{none}
\newcommand{\BookEndAppBar}{45}
\newcommand{\BookEndBarVerdict}{no bar}
\newcommand{\ConversionPlateau}{$+$13.3}
\newcommand{\ConversionPlateauLo}{$+$11.5}
\newcommand{\ConversionPlateauHi}{$+$15.1}
\newcommand{\ConversionBar}{none}
\newcommand{\ConversionAppBar}{30}
\newcommand{\ConversionBarVerdict}{no bar}
\newcommand{\ForcingPlateau}{$+$8.1}
\newcommand{\ForcingPlateauLo}{$+$7.0}
\newcommand{\ForcingPlateauHi}{$+$9.2}
\newcommand{\ForcingBar}{none}
\newcommand{\ForcingAppBar}{30}
\newcommand{\ForcingBarVerdict}{no bar}
\newcommand{\StructurePlateau}{$-$0.5}
\newcommand{\StructurePlateauLo}{$-$6.0}
\newcommand{\StructurePlateauHi}{$+$5.5}
\newcommand{\StructureBar}{none}
\newcommand{\StructureAppBar}{30}
\newcommand{\StructureBarVerdict}{no bar}
\newcommand{\Bars}{zero}
\newcommand{\TrendTable}{%
\begin{tabular}{llllll}
\toprule
mover's rating & none & book end & conversion & forcing & structure \\
\midrule
900--1099 & $+$5.3 {\footnotesize ($+$4.2, $+$6.4)} & $+$8.8 {\footnotesize ($+$4.3, $+$14.2)} & $+$10.0 {\footnotesize ($+$8.9, $+$11.1)} & $+$7.8 {\footnotesize ($+$6.9, $+$8.7)} & $+$1.1 {\footnotesize ($-$5.6, $+$7.6)} \\
1100--1299 & $+$4.9 {\footnotesize ($+$4.2, $+$5.6)} & $+$2.8 {\footnotesize ($+$0.7, $+$5.0)} & $+$8.7 {\footnotesize ($+$7.9, $+$9.4)} & $+$7.6 {\footnotesize ($+$7.0, $+$8.2)} & $+$3.9 {\footnotesize ($-$0.4, $+$8.3)} \\
1300--1499 & $+$4.6 {\footnotesize ($+$4.0, $+$5.1)} & $+$2.2 {\footnotesize ($+$0.3, $+$4.2)} & $+$8.9 {\footnotesize ($+$8.2, $+$9.6)} & $+$7.7 {\footnotesize ($+$7.2, $+$8.1)} & $-$0.1 {\footnotesize ($-$3.0, $+$3.1)} \\
1500--1699 & $+$4.5 {\footnotesize ($+$3.9, $+$5.1)} & $+$1.2 {\footnotesize ($-$0.7, $+$3.3)} & $+$8.8 {\footnotesize ($+$8.1, $+$9.6)} & $+$8.2 {\footnotesize ($+$7.6, $+$8.7)} & $+$3.8 {\footnotesize ($+$0.3, $+$7.5)} \\
1700--1899 & $+$5.0 {\footnotesize ($+$4.3, $+$5.7)} & $+$1.3 {\footnotesize ($-$0.8, $+$3.8)} & $+$8.9 {\footnotesize ($+$7.9, $+$10.0)} & $+$7.1 {\footnotesize ($+$6.4, $+$7.8)} & $-$3.0 {\footnotesize ($-$6.0, $+$0.5)} \\
1900--2099 & $+$4.0 {\footnotesize ($+$3.4, $+$4.6)} & $+$0.8 {\footnotesize ($-$1.0, $+$2.9)} & $+$8.8 {\footnotesize ($+$7.7, $+$9.9)} & $+$7.4 {\footnotesize ($+$6.8, $+$8.1)} & $+$5.8 {\footnotesize ($+$1.4, $+$10.7)} \\
\bottomrule
\end{tabular}
}
\newcommand{\SamePositions}{292}
\newcommand{\SameQuick}{9,974}
\newcommand{\SameLong}{386}
\newcommand{\SameDelta}{$+$1.9}
\newcommand{\SameLo}{$-$1.0}
\newcommand{\SameHi}{$+$5.0}
\newcommand{\SameCutBound}{1.0}
\newcommand{\SameBoards}{%
\begin{subfigure}[t]{0.32\textwidth}\centering
\chessboard[setfen={rnbqkbnr/ppp1pppp/8/3p4/3P1B2/8/PPP1PPPP/RN1QKBNR b KQkq - 0 1}, boardfontsize=11pt]
\caption*{Black to move: 655 quick moves, 0.0\% serious; 6 long, 0.0\%.\\[0.2em]\fen{rnbqkbnr/ppp1pppp/8/3p4/3P1B2/8/PPP1PPPP/RN1QKBNR b KQkq -}}
\end{subfigure}\hfill
\begin{subfigure}[t]{0.32\textwidth}\centering
\chessboard[setfen={r1bqkbnr/pppp1ppp/2n5/4p3/2B1P3/5N2/PPPP1PPP/RNBQK2R b KQkq - 0 1}, boardfontsize=11pt]
\caption*{Black to move: 477 quick moves, 6.9\% serious; 4 long, 0.0\%.\\[0.2em]\fen{r1bqkbnr/pppp1ppp/2n5/4p3/2B1P3/5N2/PPPP1PPP/RNBQK2R b KQkq -}}
\end{subfigure}\hfill
\begin{subfigure}[t]{0.32\textwidth}\centering
\chessboard[setfen={rnbqkbnr/pp1ppppp/8/2p5/4P3/8/PPPP1PPP/RNBQKBNR w KQkq - 0 1}, boardfontsize=11pt]
\caption*{White to move: 324 quick moves, 0.0\% serious; 1 long, 0.0\%.\\[0.2em]\fen{rnbqkbnr/pp1ppppp/8/2p5/4P3/8/PPPP1PPP/RNBQKBNR w KQkq -}}
\end{subfigure}
}
\newcommand{\SensitivityTable}{%
\begin{tabular}{lllll}
\toprule
 & \multicolumn{2}{c}{ten strata} & \multicolumn{2}{c}{quick from 6 seconds} \\
\cmidrule(lr){2-3}\cmidrule(lr){4-5}
moment & $\Delta$ & $\Delta - \Delta_{\text{none}}$ & $\Delta$ & $\Delta - \Delta_{\text{none}}$ \\
\midrule
Book end & $+$1.2 {\footnotesize ($-$0.4, $+$3.0)} & $-$3.6 {\footnotesize ($-$5.3, $-$1.8)} & $+$0.5 {\footnotesize ($-$1.7, $+$2.8)} & $-$3.5 {\footnotesize ($-$5.8, $-$1.2)} \\
Conversion & $+$8.8 {\footnotesize ($+$8.1, $+$9.4)} & $+$3.9 {\footnotesize ($+$3.1, $+$4.8)} & $+$7.2 {\footnotesize ($+$6.5, $+$8.0)} & $+$3.2 {\footnotesize ($+$2.3, $+$4.2)} \\
Forcing & $+$7.4 {\footnotesize ($+$6.9, $+$7.9)} & $+$2.6 {\footnotesize ($+$1.9, $+$3.3)} & $+$6.4 {\footnotesize ($+$5.9, $+$7.0)} & $+$2.4 {\footnotesize ($+$1.6, $+$3.2)} \\
Structure & $+$0.8 {\footnotesize ($-$2.2, $+$3.9)} & $-$4.0 {\footnotesize ($-$7.0, $-$0.9)} & $-$1.2 {\footnotesize ($-$4.8, $+$2.7)} & $-$5.2 {\footnotesize ($-$8.9, $-$1.3)} \\
\bottomrule
\end{tabular}
}
\newcommand{\SeriousRate}{14.5}
\newcommand{\ConcentrationTable}{%
\begin{tabular}{lrrrr}
\toprule
moment & moves & share of moves & share of serious errors & ratio \\
\midrule
book end & 38,411 & 3.9 & 1.6 & 0.40 \\
conversion & 209,836 & 21.4 & 21.2 & 0.99 \\
forcing & 445,443 & 45.4 & 49.9 & 1.10 \\
structure & 6,929 & 0.7 & 0.5 & 0.76 \\
none & 280,327 & 28.6 & 26.8 & 0.94 \\
\midrule
in check & 61,508 & 6.3 & 4.4 & 0.70 \\
after a capture & 253,356 & 25.8 & 22.5 & 0.87 \\
his piece hanging & 375,658 & 38.3 & 44.1 & 1.15 \\
opponent's piece hanging & 252,969 & 25.8 & 26.6 & 1.03 \\
after the opponent's error & 133,873 & 13.6 & 33.1 & 2.43 \\
defending, $-$3 or worse & 178,114 & 18.2 & 6.0 & 0.33 \\
book end and his next four & 193,286 & 19.7 & 12.4 & 0.63 \\
high danger, top fifth & 227,445 & 23.2 & 54.8 & 2.37 \\
\bottomrule
\end{tabular}
}
\newcommand{\AskedSeconds}{541}
\newcommand{\AskedShare}{90}
\newcommand{\HisGames}{632}
\newcommand{\HisMoves}{16,290}
\newcommand{\HisTable}{%
\begin{tabular}{lrrrrrr}
\toprule
 & \multicolumn{2}{c}{per game} & \multicolumn{2}{c}{share of long thinks} & \multicolumn{2}{c}{his serious errors} \\
\cmidrule(lr){2-3}\cmidrule(lr){4-5}\cmidrule(lr){6-7}
moment & his & his level & his & his level & quick {\footnotesize ($n$)} & long {\footnotesize ($n$)} \\
\midrule
book end & 0.9 & 1.0 & 1.4 & 1.1 & 1.5 {\footnotesize (525)} & 0.0 {\footnotesize (8)} \\
conversion & 4.8 & 5.3 & 8.3 & 6.3 & 8.6 {\footnotesize (1,735)} & 18.2 {\footnotesize (253)} \\
forcing & 12.2 & 11.2 & 6.1 & 5.5 & 9.7 {\footnotesize (5,040)} & 23.7 {\footnotesize (473)} \\
structure & 0.2 & 0.2 & 3.8 & 4.4 & 10.3 {\footnotesize (107)} & 33.3 {\footnotesize (6)} \\
none & 7.6 & 7.0 & 5.7 & 5.6 & 6.9 {\footnotesize (3,163)} & 16.7 {\footnotesize (270)} \\
\bottomrule
\end{tabular}
}
\newcommand{\HisMinLong}{50}
\newcommand{\HisConversionLong}{8.3}
\newcommand{\PeersConversionLong}{6.3}
\newcommand{\AfterMoves}{14}
\newcommand{\AfterErrors}{33}
\newcommand{\HighMoves}{23}
\newcommand{\HighErrors}{55}
\newcommand{\BookEndMoves}{3.9}
\newcommand{\BookEndErrors}{1.6}
\newcommand{\TriggersPerGame}{18}
\newcommand{\AnalysedPerGame}{25}
\newcommand{\BookEndErrorsPerGame}{0.06}
\newcommand{\AfterErrorsPerGame}{1.2}
\newcommand{\ErrorsPerGame}{3.6}
\newcommand{\TriggerShare}{71}
\newcommand{\ForcingShare}{45}


\title{Where does a longer think pay in rapid?\\[0.3em]
\large The app's long-think triggers on \GamesBuilt{} Lichess 10-minute games}
\author{Chess Coach research lab}
\date{}

\begin{document}
\maketitle

\begin{abstract}
The app tells the player to stop and think at four kinds of moment: the first move out of the
opening book (\BookEnd{} seconds), a winning position to convert, a forcing moment, and a pawn
capture outside the endgame (\Deep{} seconds each). On \GamesBuilt{} rated 10-minute Lichess games
from \Month, built with the app's own code, players at his level made serious errors on
\NoneRateQuick\% of ordinary moves played in under \Rushed{} seconds and on \NoneRateLong\% of those
played after \Deep{} seconds or more. Compared only in positions a danger score rates alike, a long
think was followed by more serious errors at ordinary moves (\NoneDelta{} points), at conversion
moments (\ConversionDelta) and at forcing moments (\ForcingDelta), and by no fewer at the book exit
(\BookEndDelta) or a pawn capture (\StructureDelta). By the rule fixed before any effect was
computed, no trigger shows that a long think pays, none of eight other moments does, and no think
length is a bar. The pattern holds at every rating from \RatingLo{} to \RatingHi{} and, at the book exit, within
the very same position. Players think long where a position is harder than the board and the
clocks show, and that outweighs whatever the think buys: this design cannot see past it. It does
show two things the app can use. The triggers fire on \TriggerShare\% of the moves analysed and ask
for about \AskedSeconds{} of a game's 600 seconds, so they cannot be followed as written. And
serious errors gather where the triggers do not point: a third of them come on the move right
after the opponent's own serious error, while the book exit holds \BookEndErrors\% of them.
\end{abstract}

\section{The question}

The app's method service (\texttt{triggerOf}, at commit \texttt{\AppCommit}) names the moments
where a long think should pay. The first move out of the player's book should get \BookEnd{}
seconds. A position to convert (an evaluation of $+3$ or better), a forcing moment (in check, just
after a capture, a piece hanging on either side, or a move that gives check) and a pawn capture
outside the endgame should get \Deep. The app calls a move made in under \Rushed{} seconds rushed
and one of \Deep{} seconds or more a long think. For players at his level in 10-minute games, does
a long think cut serious errors at these moments, and more than at an ordinary move? Do other
moments pay more? Are \Deep{} and \BookEnd{} seconds the right bars?

\section{The data}

The Lichess database for \Month{} (CC0) holds \GamesRead{} games, \GamesRapid{} of them rated
rapid; \GamesWithEvals{} of those were played at 10 minutes with no increment and carry Lichess's
engine evaluation on every move beside the clock. The project's stream kept, in file order, the
games that ended normally or on time, had at least 20 plies and two ratings within 300 points,
by bins of the mean rating (Table~\ref{tab:kept}).

\begin{table}[h]
\centering
\small\KeptTable
\caption{Games kept, by the players' mean rating.}
\label{tab:kept}
\end{table}

Each game was read as the routine's archive row and built with the app's own code through the
lab's bridge: the metrics service's \texttt{primaryFacts} and \texttt{deriveGameFacts} with the
routine's ports, and \texttt{triggerOf}. Every game was built once with each player as ``me'', so
every move gets its moment from its own side; Black's build is White's with the evaluations and
the result turned around, and on \BlackChecked{} games it gave every move the same facts as a real
Black build.
All \GamesBuilt{} games built. Two choices are this study's. A Lichess player has no repertoire in
the app, so the book exit is the first move the app's review reader does not call book (its book
covers the first 12 plies). And a serious error is a mistake or a blunder as the reader classifies
it from the evaluations; the app also counts misses, which need the routine's miss ledger.

The moves studied have a think time (each side's first move has none), are graded and not trivial
(the app's \texttt{nonTrivial}: not book, not forced, a win probability before the move between 5
and 95\%), and leave the mover at least \MinClock{} seconds on the clock before the move, so a long
think was affordable: \MovesAnalysed{} of \MovesTimed{} timed moves. His chess.com rapid rating
comes to about \WindowCentre{} on Lichess by the app's default conversion, so the rule is applied to
the moves of players rated \WindowLo{} to \WindowHi: \MovesPrimary{} moves in \PlayerGames{} player-games, of
which \QuickShare\% were quick and \LongShare\% long.

\section{The method}

Players think longer where a position is harder, and harder positions breed errors, so long and
quick moves are compared only in positions a \emph{danger score} rates alike. The score is the
chance that a quick move in the position is a serious error, predicted from what the player could
see before moving: both ratings and clocks, the opponent's last think, the ply and the moves since
the book exit, the evaluation and win probability before the move, the phase, the material, the
legal moves, the checks and captures available, the hanging pieces on each side, being in check,
coming just after a capture or either side's serious error, and the colour. Nothing about the
think or the move played. A gradient-boosted model (LightGBM, \Trees{} trees, fixed settings) was
trained on the quick moves of every rating and cross-fitted in two halves split by game, so each
move's score comes from a model that never saw its game. It ranks quick moves well (AUC \AucQuick)
and the others nearly as well (\AucOther).

For a set of moves $T$, its moves are split into \Strata{} strata at the quintiles of the danger
score over $T$'s quick moves, and $\Delta_T$ is the long moves' serious-error rate minus the quick
moves', averaged over the strata with weights equal to $T$'s quick moves in each: how much a long
think would change the error rate of the moves now played quickly. Intervals come from \Draws{}
bootstrap draws over games (Poisson weights, seed \Seed), the 95\% percentile interval, and a
$p$-value of twice the smaller share of draws on either side of zero; every comparison uses the
same draws.

The rule, committed before any effect was computed: a long think \emph{pays} at $T$ when
$\Delta_T$'s interval lies below zero; $T$ \emph{marks} where it pays when it also beats ordinary
moves ($\Delta_T - \Delta_{\text{none}}$ below zero with a Holm-adjusted $p$ under 0.05 over the
four triggers). Each trigger is kept when both hold. Eight candidate moments, each defined before
the move, would join the triggers under the same test over eight. A trigger's bar is the shortest
think bin whose $\Delta$ reaches \BarShare\% of a \Plateau-second think's, when that one pays.

\section{Results}

\subsection{The triggers}

\begin{table}[h]
\centering
\small\setlength{\tabcolsep}{4pt}\TriggerTable
\caption{The four triggers and ordinary moves (none), players rated \WindowLo--\WindowHi.
Per game: moves per player per game among those analysed. Quick: under \Rushed{} seconds; long:
\Deep{} seconds or more. Error rates are raw; $\Delta$ is long minus quick within danger strata,
in percentage points, with its 95\% interval.}
\label{tab:triggers}
\end{table}

At ordinary moves a long think is followed by \NoneDelta{} points more serious errors than a quick
move in a position the score rates alike (Table~\ref{tab:triggers}). At conversion and forcing
moments the excess is larger still (\ConversionDelta{} and \ForcingDelta), and larger than at
ordinary moves by \ConversionVsNone{} and \ForcingVsNone{} points. At the book exit and at a pawn
capture, the intervals hold zero (\BookEndLo{} to \BookEndHi{} and \StructureLo{} to
\StructureHi), and the excess is smaller than at ordinary moves (by \BookEndVsNone{} and
\StructureVsNone{} points, Holm $p$ \BookEndPHolm{} and \StructurePHolm). No trigger's interval
lies below zero, so the verdict is the same for all four: \emph{no sign it pays}. Read as the rule
reports it, a long think's extra minute at a conversion moment goes with \ConversionSaved{}
serious errors saved (\ConversionSavedLo{} to \ConversionSavedHi), at a forcing moment with
\ForcingSaved{}: a cost, not a saving.

\subsection{Eight other moments}

\begin{table}[h]
\centering
\footnotesize\setlength{\tabcolsep}{4pt}\CandidateTable
\caption{The candidates, each defined by what the player sees before moving. Joins: pays, and
beats ordinary moves with a Holm-adjusted $p$ under 0.05 over the eight.}
\label{tab:candidates}
\end{table}

None pays (Table~\ref{tab:candidates}): at every candidate a long think is followed by more serious
errors than a quick one, most of all right after the opponent's serious error
(\CandAfterOppErrorDelta) and in the top fifth of the danger score (\CandHighDangerDelta). The
first four candidates are the parts of the forcing trigger; none of them carries it.

\subsection{The bars}

\begin{figure}[h]
\centering
\begin{subfigure}{0.49\textwidth}\centering\BookEndBarPlot\end{subfigure}\hfill
\begin{subfigure}{0.49\textwidth}\centering\ConversionBarPlot\end{subfigure}

\medskip
\begin{subfigure}{0.49\textwidth}\centering\ForcingBarPlot\end{subfigure}\hfill
\begin{subfigure}{0.49\textwidth}\centering\StructureBarPlot\end{subfigure}
\caption{$\Delta$ for each length of think against the quick moves, by trigger, in percentage
points with 95\% intervals. Below zero would mean fewer serious errors after the longer think.}
\label{fig:bars}
\end{figure}

No think length pays at any trigger, so no trigger has a bar (Figure~\ref{fig:bars}). At
conversion and forcing moments the excess grows with the length of the think: every bin from
\Deep{} seconds up lies above the 10--14-second bin. The longer the think, the harder the position
it marks, in ways the danger score does not see.

\subsection{Every rating}

\begin{table}[h]
\centering
\footnotesize\setlength{\tabcolsep}{2pt}\TrendTable
\caption{$\Delta$ by the mover's rating, in percentage points with 95\% intervals. Described,
not tested.}
\label{tab:trend}
\end{table}

The answer does not move with the level (Table~\ref{tab:trend}): at every rating from \RatingLo{} to \RatingHi,
ordinary, conversion and forcing moves are followed by more serious errors after a long think. So
the uncertainty in his Lichess level matters little here.

\subsection{The same position}

A book-exit move is often made in a position many games reach, so long and quick moves can be
compared in the very same position, with no danger score at all: \SamePositions{} positions were
reached with both (\SameQuick{} quick moves, \SameLong{} long). Weighted by their quick moves, the
long thinks were followed by \SameDelta{} points more serious errors (\SameLo{} to \SameHi): a cut
of more than \SameCutBound{} points is ruled out. The positions most games reached are
quiet (Figure~\ref{fig:same}): in them long thinks are rare, and serious errors rarer than at his
level's average move.

\begin{figure}[h]
\centering
\SameBoards
\caption{The three book-exit positions the most games at his level reached, with the moves made
in them.}
\label{fig:same}
\end{figure}

\subsection{Sensitivity}

With ten strata instead of five, and with quick moves limited to those over \Decided{} seconds (the
app's own clock buckets treat faster moves as decided before the turn), no trigger pays and each
comparison with ordinary moves keeps its side (Table~\ref{tab:sens}).

\begin{table}[h]
\centering
\small\setlength{\tabcolsep}{4pt}\SensitivityTable
\caption{The triggers' $\Delta$ and $\Delta - \Delta_{\text{none}}$ under two other choices.}
\label{tab:sens}
\end{table}

\section{Where the errors are}

This section was added after the rule was committed, before the month's results were computed, and
tests nothing. A trigger can save at most the errors made at its moment, so it is worth knowing
where they are. At his level a player makes about \ErrorsPerGame{} serious errors a game among the
moves analysed (\SeriousRate\% of them). Table~\ref{tab:where} sets each moment's share of the
moves against its share of those errors.

\begin{table}[h]
\centering
\small\ConcentrationTable
\caption{Each moment's share of the moves analysed at his level and of their serious errors, at
every length of think. Ratio: the second over the first. The candidates overlap.}
\label{tab:where}
\end{table}

The right moment to stop is where the errors are. The move right after the opponent's serious error
is \AfterMoves\% of the moves and \AfterErrors\% of the serious errors, about \AfterErrorsPerGame{}
a game. By the evaluations, failing to take what an error offers costs as much as making one, so
these include every punishment missed. The top fifth of the danger score holds \HighErrors\%. The book exit, where the app asks
for its longest think, is \BookEndMoves\% of the moves and \BookEndErrors\% of the errors, about
\BookEndErrorsPerGame{} a game: even a think that removed every one of them would save little.
Meanwhile the triggers fire about \TriggersPerGame{} times in a player's game, on \TriggerShare\% of
the \AnalysedPerGame{} or so moves analysed (the forcing trigger alone on \ForcingShare\%), and at
the app's bars they ask for about \AskedSeconds{} seconds, \AskedShare\% of a 10-minute game.

\section{His games}

\begin{table}[h]
\centering
\small\setlength{\tabcolsep}{4pt}\HisTable
\caption{His \HisGames{} 10-minute games in the app's archive, read the same way and only
described. His error rates come from the routine's engine rather than Lichess's, so they are not
set beside his level's.}
\label{tab:his}
\end{table}

His games meet the moments about as often as his level does (Table~\ref{tab:his}), except the
rare pawn capture outside the endgame, which comes up more often in his. He thinks long
on \HisConversionLong\% of his conversion moves, against \PeersConversionLong\% at his level; at
the other moments his share is close to theirs. After a long think his own serious errors run
above those after a quick move at every moment with at least \HisMinLong{} long thinks to read, as at
his level.

\section{What it means for the app}

\begin{enumerate}[nosep]
\item \textbf{The triggers' think times are advice, not a measurement.} Nothing here shows a long
think cutting serious errors at any moment, for any rating. That does not show the advice is wrong:
it shows that where players choose to think long, the position is harder than anything the board
and the clocks reveal, by more than the think can make up. The app should not tell him a think
time at a moment saves errors as if it had been measured.
\item \textbf{Fewer triggers.} As written they fire on most of the moves he plays and ask for
nine tenths of his clock. A trigger list a player can follow in 10 minutes has to be short, and
the forcing trigger, which fires on almost half the moves, is the place to cut.
\item \textbf{Point them where errors are.} The move after the opponent's serious error holds a
third of the serious errors at his level on a seventh of the moves; it is also the ledger's own
``after a gift'' leak. The book exit holds almost none: its \BookEnd{} seconds are the cheapest to
give back.
\item \textbf{Measure the effect with players, not positions.} Whether a long think pays at a
moment can be read by comparing players who habitually think longer there with others of their
rating, in their other games, against the same comparison at ordinary moves; or, on his own games,
by having the coach flag a moment in some games and not in others.
\end{enumerate}

\section{Limits}

The study is observational. The danger score adjusts for what the board and the clocks show, not
for what the player saw; the comparison with ordinary moves is less exposed when the bias is alike,
but it is not alike at every moment. Two parts of the triggers depend on the move played (giving
check, a pawn capture), which a long think may change. The clock a long think takes from later
moves is not counted. The book exit is the reader's 12-ply book's, not a repertoire's, and a serious
error leaves out misses. Lichess evaluations come from its own server analysis; games someone asked
to have analysed are not a random sample; and his Lichess level is the app's conversion of his
chess.com rating, not a measurement, though the answer is the same at every level.

\end{document}
